\section{Attachment maps, monodromy and the spectral hierarchy}
\label{sec:topology-comparison}

\subsection{Relations among vanishing cycles}

Let $Y$ be the small resolution, $U$ the common exterior,
$C_i$ the exceptional curves, and $\delta_i$ the oriented
vanishing spheres. The integral Thom and Mayer--Vietoris
calculation in \cite[Section~3.3, proposition on restriction
and spectral coefficient]{wuebben-hodge} gives
\begin{equation}\label{eq:integral-conifold-sequence}
 \begin{aligned}
 0\longrightarrow H^2(X_p;\ZZ)
 &\longrightarrow H^2(Y;\ZZ)
 \xrightarrow{\ r\ }\ZZ^8
 \xrightarrow{\ \partial\ }H^3(X_p;\ZZ)\\
 &\longrightarrow H^3(U;\ZZ)\longrightarrow0,
 \end{aligned}
\end{equation}
where $r(D)_i=D\cdot C_i$ and
$\partial e_i=\operatorname{PD}_L(\delta_i)$.
The latter class is defined by the oriented Thom
construction, including any torsion in its global image.

Over $\QQ$, put
$Q_{\mathrm{lost}}=H^2(Y;\QQ)/\im H^2(X_p;\QQ)$
and $V_{\mathrm{van}}=\im\partial$. Then
\begin{equation}\label{eq:rational-conifold-sequence}
 0\longrightarrow Q_{\mathrm{lost}}
 \xrightarrow{\ r\ }\QQ^8
 \xrightarrow{\ \partial\ }V_{\mathrm{van}}
 \longrightarrow0.
\end{equation}
Their dimensions are two and six, respectively.
In our quotient basis the columns of $r$ are
$(1,1,1,1;0,0,0,0)$ and $(0,0,0,0;1,1,1,1)$.
Thus the lost degree-two cohomology quotient is the space of
relations among the cohomological vanishing cycles,
and its composite to $V_{\mathrm{van}}$ is zero.
The exceptional-curve map $s:\QQ^8\to H_2(Y;\QQ)$,
$s(e_i)=[C_i]$, is dual to $r$.
Its six-dimensional relation space $\ker s$ is
a different space, paired with the period
coordinates of Section~\ref{sec:geometry}.

Lee--Lin--Wang realize
\eqref{eq:rational-conifold-sequence} as the
rational weight-two basic exact sequence of a
semistable model \cite[Theorem~1.14]{lee-lin-wang};
their Proposition~3.16 treats vanishing-period
coordinates. These are classical deformation
and Hodge-theoretic inputs. Rational exactness
does not by itself establish integral saturation
or recover an unknown attaching map.

\subsection{The nearby-cycle sequence}

For a specified proper one-parameter
ordinary-double-point smoothing $f$ with
smooth total space, the unshifted nearby-
and vanishing-cycle triangle is
\[
 i^*\QQ\longrightarrow R\psi_f\QQ
 \longrightarrow R\phi_f\QQ
 \longrightarrow i^*\QQ[1].
\]
The Milnor fibres of the threefold nodes
give the specialization segment
\begin{equation}\label{eq:nearby-cycle-sequence}
 \begin{aligned}
 0\longrightarrow H^3(X_0;\QQ)
 &\longrightarrow H^3(X_t;\QQ)
 \xrightarrow{\ \mathrm{can}\ }\QQ^8\\
 &\longrightarrow H^4(X_0;\QQ)
 \longrightarrow H^4(X_t;\QQ)\longrightarrow0.
 \end{aligned}
\end{equation}
The map $\mathrm{can}$ evaluates on vanishing
spheres. Its image annihilates the relations
in \eqref{eq:rational-conifold-sequence}.
The map $\partial$ there is its
Poincaré--excision dual, with the chosen
orientation convention; it is not the arrow
from $\QQ^8$ to $H^4(X_0;\QQ)$ in
\eqref{eq:nearby-cycle-sequence}.
For the triangles, stalks and canonical/variation
identities see \cite[Section~3]{massey}, and
for the conifold specialization see
\cite[Section~1]{lee-lin-wang}.
After ramified base change the total space
need not be smooth; one must use nearby-cycle
base change or a semistable model instead of
applying the smooth-total-space local argument
without its hypothesis.

\subsection{What monodromy does not determine}

Picard--Lefschetz gives
$N_i\gamma=\pm\langle\gamma,\delta_i\rangle\delta_i$.
The vanishing spheres are disjoint, so
$N_iN_j=0$. For an all-node loop with
positive winding numbers $\nu_i$,
\begin{equation}\label{eq:all-node-monodromy}
 T=I+N,\qquad N=\sum_i\nu_iN_i,\qquad N^2=0
\end{equation}
as in \cite[Section~3.3, equation (3.4)]{lee-lin-wang}.
The monodromy image is abelian; no assertion
that the fundamental group of the
discriminant complement is abelian is needed.

The weight filtration centred at three is
\begin{equation}\label{eq:monodromy-weight-filtration}
 0=W_1\subset W_2=\im N
       \subset W_3=\ker N\subset W_4=H^3.
\end{equation}
Here $\dim W_2=6$. These rational subspaces
are independent of the positive $\nu_i$.
In a symplectic basis let $B$ have columns
$\delta_i$ and let $J$ be the intersection
matrix. Then $N=\pm BDB^{\mathsf T}J$,
$D=\operatorname{diag}(\nu_i)>0$.
The matrix $BDB^{\mathsf T}$ has the
same rank and image as $B$, so
$\im N=\im B$ and
$\ker N=(\im B)^{\perp_J}$.
Integral image indices can change under
ramification. The characterization of the
weight filtration is
\cite[Section~6, Lemma~(6.4)]{schmid};
the identification of $V_{\mathrm{van}}$
with $W_2$ is
\cite[Lemma~1.12]{lee-lin-wang}.

The two admissible paths
\[
 (a,b)=(e^{-h^{-1}},e^{-h^{-2}})
 \quad\hbox{and}\quad
 (a,b)=(e^{-h^{-1}},e^{-h^{-3}})
\]
have the same marked central attachment
sequence and local monodromy representation.
Choose the same reference all-node loop;
its operator and weight filtration are
then the same. Nevertheless
\eqref{eq:actual-general-hierarchy} gives
the actual spectral orders $(h,h^2)$ and
$(h,h^3)$. The stated topology and monodromy
data therefore do not determine the
capacity hierarchy.

An interval path has no distinguished loop,
and a nonanalytic path does not define an
ordinary holomorphic-disc nearby-cycle
functor. We make no assertion that the
full limiting mixed Hodge structures or
period asymptotics agree on these paths.
The limiting Hodge filtration contains
further information; its coordinate
dependence by $\exp(cN)$ is recorded
before \cite[Theorem~(6.16)]{schmid}.

\subsection{Comparison with adiabatic spectral sequences}

For a fixed compact fibration,
Mazzeo--Melrose identify the adiabatic
Hodge problem with the Leray spectral
sequence \cite[Theorems~(7), (15) and
Corollary~(18)]{mazzeo-melrose}.
For a fixed smooth Riemannian foliation
on a closed manifold, Álvarez
López--Kordyukov study
$g_h=h^{-2}g_\perp+g_F$.
The transverse-degree filtration gives
\[
 \dd_h=\dd_{0,1}+h\dd_{1,0}+h^2\dd_{2,-1}.
\]
Their Theorem~A counts eigenvalues,
including zeros, of order $O(h^{2k})$
by the degreewise dimension of the
$k$-th page for $k\ge2$; at $k=1$
it uses the Hausdorff quotient of
the first page \cite{alvarez-kordyukov}.
The corrected Theorem~B gives only
weak subsequential convergence at
$k=1$ in its stated setting, with
strong convergence for $k\ge2$
under its energy hypotheses.

Those results include a filtered
differential complex and its page
differentials. The successive kernels
in Theorem~\ref{thm:fixed-map-hierarchy}
come instead from the fixed restriction
rows of a finite energy form.
We have not constructed page differentials
or a compatible de Rham Laplacian in
every degree. The conifold smoothing
also changes topology, unlike the
fixed-foliation adiabatic metric.
Thus neither the capacity groups nor
the relation space in
\eqref{eq:rational-conifold-sequence}
are identified with those pages.
